% Benvenuti nel template per le relazioni di Sperimentazioni di Fisica 1 del corso triennale in Astronomia dell'Università di Padova.t
% Per compilarlo, cliccate su Menu, selezionate come compilatore "pdfLaTeX", TeX Live Version 2020 (o uno degli altri, whatever works), e come documento principale "Relazione.tex"

%%% Preambolo del documento %%%
% Non toccate questa roba, a meno che non siate abbastanza sgamati da sapere quello che state facendo, tipo vi piacciono di più certi pacchetti rispetto ad altri, o magari ve ne manca qualcuno.
\documentclass[a4paper, 12pt, twoside]{article}
\usepackage[top=2.5cm, bottom=2.5cm, inner=2.5cm, outer=2.5cm]{geometry}
% da sostituire in caso di stampa fronte-retro
%\usepackage[top=2.5cm, bottom=2.5cm, inner=4cm, outer=4cm, right=2.5cm]{geometry}
\usepackage[utf8]{inputenc} % codifica UTF-8
\usepackage[italian]{babel}
\usepackage{amsmath}
\usepackage{listings}
\usepackage{multicol}
\usepackage{textcomp}
\usepackage{multirow}
\usepackage{siunitx}
\usepackage{multicol}
\usepackage{booktabs}
\usepackage{fixmath}
\usepackage{graphicx}
\usepackage{floatflt}
\usepackage[export]{adjustbox}
\usepackage{epsfig}
% \usepackage{pstricks}
\usepackage[labelfont=bf, font=scriptsize]{caption}
\usepackage[italian]{varioref}
\usepackage[suftesi]{frontespizio}
\usepackage{color}
\usepackage{tikz}
\usepackage{caption}
\usepackage{subcaption}
\usepackage{pgfplots}
\usepackage{comment}
\usepackage{lipsum}
\usepackage{wrapfig}
\pgfplotsset{compat=1.16}
\usepackage[colorlinks=true, urlcolor=blue, linkcolor=black]{hyperref}
\usepackage{svg}
\usepackage{siunitx}
\usepackage{tabularx}

% ROBA AGGIUNTA DA SAS - se si rompe tutto eliminare questa parte
\usepackage{setspace}
\onehalfspacing
\usepackage{changepage}
\usepackage{colortbl}
\usepackage{titlesec}
\titlespacing*{\section}{0pt}{10ex plus 3ex minus 2ex}{4.3ex plus .22ex}
\titlespacing*{\subsection}{0pt}{5.4ex plus .3ex minus .2ex}{2.8ex plus .2ex}
\usepackage[italian]{babel}
\usepackage[svgnames]{xcolor}
\usepackage[outputdir=build]{minted}

\usemintedstyle{vs}


% bibliografia
\usepackage[backend=biber, style=numeric-comp, sorting=none]{biblatex}
\addbibresource{bibliografia.bib}
\AtEveryBibitem{\clearfield{urlyear}\clearfield{urlmonth}\clearfield{urlday}}

% FINE ROBA AGGIUNTA DA SAS


\begin{document}

\begin{frontespizio}

\Universita{Milano-Bicocca}
\Logo{immagini/logoUNIMIB.png}
\Dipartimento{Fisica}
\Corso[Laurea Triennale]{Fisica} 
\Annoaccademico{2025-2026}
\Titoletto{Tesi di Laurea Triennale}
\Titolo{Metodo dei fasori applicato all'analisi dei tempi di vita di fluorescenza}
\Sottotitolo{Novembre 2026}

\Candidato[898437]{Di Lernia Sara, \textsf {s.dilernia1@campus.unimib.it}}
\NRelatore{Relatrice}{} 
\Relatore{Prof.ssa Laura Sironi}
\NCorrelatore{Correlatore}{}
\Correlatore{Dott. Davide Panzeri}

\end{frontespizio}


\IfFileExists{\jobname-frn.pdf}{}{%
\immediate\write18{lualatex -interaction=nonstopmode -output-directory=build \jobname-frn.tex}}

% --- INIZIO PAGINA BIANCA ---
\newpage
\thispagestyle{empty} % Rimuove il numero di pagina
\mbox{}               % Inserisce un box vuoto invisibile
\newpage
% --- FINE PAGINA BIANCA ---

%%% Indice %%%
\vspace*{\fill}
\tableofcontents
\vspace*{\fill}
\newpage


%\input{SEZIONI/01_OBIETTIVI.tex}

\chapter{Introduzione}\label{sec:introduzione}


\begin{itemize}

    \item cosa vuol dire studiare i tempi di decadimento della fluorescenza
    \item dove fallisce il fit a esponenziale multiplo
    \item FLIM e metodo dei fasori

\end{itemize}
\section{Fondamenti teorici}\label{sec:cenni_teorici} 

\subsection{\textit{FLIM} e analisi fasoriale}

%\cite{ranjitFitfreeAnalysisFluorescence2018}

Con il termine "tempo di vita di fluorescenza" si fa riferimento al tempo medio necessario perché una molecola allo stato eccitato emetta un fotone per tornare allo stato fondamentale. Il processo richiede in media tempi dell'ordine di grandezza del nanosecondo ($10^{-9}$s).\\
Dal punto di vista operativo, lo studio dei tempi di vita di fluorescenza in laboratorio avviene mediante l'eccitazione da parte di un laser periodico del campione e la raccolta mediante un fotorivelatore dei conseguenti decadimenti.\\


La microscopia di fluorescenza tradizionale presenta alcune criticità, prima tra tutte la difficoltà nella stima della concentrazione di fluoroforo nelle regioni di tessuto biologico \cite[Sez. 22]{mastersBookReviewPrinciples2008a}. Queste infatti dipendono da fattori complessi come determinate affinità tra fluoroforo e sostanze presenti nel campione, e non esiste una stretta correlazione tra l'intensità di fluorescenza e la concentrazione del fluoroforo \cite[Sez. 22]{mastersBookReviewPrinciples2008a}.\\

Il tempo di vita di fluorescenza, invece, offre il vantaggio di essere indipendente da quest'ultima (inoltre, è indipendente dal cammino ottico) \cite{bastiaensFluorescenceLifetimeImaging1999a}, aggirando completamente il problema.\\
Questo è l'ambito di studio della \textit{FLIM} (miscroscopia a tempo di vita di fluorescenza).\\

Un altro beneficio fondamentale è la possibilità di monitorare processi biologici che avvengono allo stato eccitato, ad esempio il trasferimento di energia per risonanza (\textit{FRET}) \cite{bastiaensFluorescenceLifetimeImaging1999a}; o l'evoluzione della concentrazione relativa di NADH (\textit{nicotinammide adenine dinucleotidee}) e FAD (\textit{flavin adenin dinucleotide}), due enzimi e fluorofori intrinseci fondamentali nella regolazione del metabolismo cellulare \cite{huComparisonPhasorAnalysis2023}.\\



%%%%%%%%%%%%%%%%%%%%%%%%%%%%%%%%%%%%%%%%%%%%%%%%%%%%%%%%%%%%%%%%%%%%%%%%%%%5

\subsection{Cenni matematici sull'analisi fasoriale}

L'implementazione pratica degli esperimenti di FLIM consiste tipicamente nell'esposizione del campione a un laser a brevi impulsi periodici, e nella successiva rilevazione del segnale di fluorescenza mediante un fotorivelatore \cite{mastersBookReviewPrinciples2008a}.\\
La probabilità che una molecola di fluoroforo si trovi allo stato eccitato al tempo $t$ segue un'andamento esponenziale decrescente, dunque il segnale di fluorescenza totale di una regione del campione contenente più molecole di un singolo fluoroforo avrà forma:

 \begin{equation*}
    I(t) = I_0\,e^{-t/\tau}
\end{equation*}

% mmmm qui c'è qualche conclusione da trarre, in più f per cosa cazzo sta


\subsubsection{Interpolazione multi-esponenziale dei decadimenti}

Quando in un campione sono presenti più fluorofori, si può immaginare che il segnale di fluorescenza totale sia una combinazione lineare di esponenziali decrescenti:

\begin{equation}
    I(t)=\sum_{j}\alpha_j\,e^{-t/\tau_j}+C \quad \quad \sum_{j} \alpha_j = 1
    \label{eq:exponential_intensity}
\end{equation}

dove $\alpha_j$ indica la frazione del singolo fluoroforo sul totale, e $\tau_j$ il suo tempo di vita; $C$ è un parametro relativo al rumore di fondo \cite{huComparisonPhasorAnalysis2023}.

Integrando l'intensità totale del segnale nel tempo (se si ignora il rumore per semplicità), si ottiene:

\begin{equation*}
    A = \int_{0}^{\infty}I(t)\,dt=\int_{0}^{\infty}\left(\sum_{j}\alpha_j\,e^{-t/\tau_j}\right)dt
\end{equation*}

Se la concentrazione di fluorofori è indipendente dal tempo (come per ogni caso presente in questa trattazione), si può portare $\alpha_j$ fuori dall'integrale:

\begin{equation*}
    A = \sum_{j} \alpha_j \left(\int_{0}^{\infty} e^{-t/\tau_j}\right) dt = \sum_{j} \alpha_j \left[ -\tau_j e^{-t/\tau_j} \right]_0^\infty = \sum_{j} \alpha_j \tau_j
\end{equation*}


Tenendo conto della normalizzazione dei rapporti di intensità definita nell'equazione \ref{eq:exponential_intensity}, si conclude: 

\begin{equation}
    \tau_m = \sum_{j}\alpha_j\,\tau_j
    \label{eq:tau_medio}
\end{equation}

Dunque, se si conosce il numero di fluorofori presenti nel campione, è possibile ricavare il loro rapporto di concentrazione mediante un \textit{fit} multi-esponenziale \cite{mastersBookReviewPrinciples2008a}.\\

Il processo di interpolazione diventa tuttavia progressivamente più complicato (e computazionalmente costoso) all'aumentare del numero di fluorofori nel campione; inoltre bisogna considerare lo svantaggio dovuto alla presenza della IRF nel segnale di intensità, che deve essere deconvoluto per ottenere il segnale di fluorescenza puro (come trattato estensivamente nella sezione succ)\cite{huComparisonPhasorAnalysis2023}.\\


%%%%%%%%%%%%%%%%%%%%%%%%%%%%%%%%%%%%%%%%%%%%%%%%%%%%%%%%%%%%55

\subsubsection{Trasformazione fasoriale}

Un metodo alternativo al fit multi-esponenziale è l'analisi con il metodo dei fasori, una tecnica che identifica ogni decadimento di fluorescenza mediante un singolo punto nel piano dei fasori, identificato dalle coordinate $S$ e $G$, definite dai primi due termini della trasformata di Fourier del segnale di fluorescenza: 

\begin{equation}
    G_{i,j}(\omega)=\frac{\int_{0}^{T}I_{i,j}\,\cos(\omega t)\,dt}{\int_{0}^{T}I_{i,j}\,dt} 
    \qquad \qquad
    S_{i,j}(\omega)=\frac{\int_{0}^{T}I_{i,j}\,\sin(\omega t)\,dt}{\int_{0}^{T}I_{i,j}\,dt}
    \label{eq:fourier_S_G_integral}
\end{equation}

dove $\omega= 2 \pi f$, $f$ frequenza del laser \cite{digmanPhasorApproachFluorescence2008}.\\

È possibile dimostrare che a un segnale periodico di fluorescenza a singolo fluoroforo (dunque mondo-esponenziale decrescente) è sempre associato un fasore sul semicerchio unitario  a $S$ positiva, e la sua posizione su di esso è determinato dal valore di $\tau$ (tempi di vita minori sulla destra).\\



% $$A = \int_0^\infty I_{sig}(t) dt = \int_0^\infty \left( \sum_j \alpha_j e^{-t/\tau_j} \right) dt$$


%__________________________________
% \begin{equation}
%     G_{i,j}(\omega)=\frac{\sum_{t=1}^{N_{bins}}I_{i,j}*\cos(\omega t)}{\sum_{t=1}^{N_{bins}}I_{i,j}} 
%     \qquad \qquad
%     S_{i,j}(\omega)=\frac{\sum_{t=1}^{N_{bins}}I_{i,j}*\sin(\omega t)}{\sum_{t=1}^{N_{bins}}I_{i,j}}
%     \label{eq:fourier_S_G}
% \end{equation}
%___________________________________


%______________________________________________________________
Se in un campione sono presenti più fluorofori (o classi di fluorofori a uguale $\tau$), si dimostra che \cite{torradoLinearCombinationProperties2022}:

\begin{equation}
    G = \sum_{k} f_k G_k, \qquad S = \sum_{k} f_k S_k \qquad \text{con} \quad \sum_{k} f_k = 1
    \label{eq:phasor_additivity}
\end{equation}


\begin{itemize}

    \item S e G piano complesso sui cui rappresento il secondo coefficiente della SdF, diviso per il primo coefficiente come normalizzazione
    \item Implementazione della teoria su in immagine divisa in pixel spaziali e bin temporali (eq \ref{eq:fourier_S_G})
    \item Fooorse riprendo/sistemo parte di teoria della rappr di decadimenti sul piano dei fasori
    \item paragrafo: cosa succede quando ho più fluorofori nello stesso punto, con intensità che diventa la somma (\ref{eq:phasor_additivity})
    \item angolo di fase e modulazione (da \cite{mastersBookReviewPrinciples2008a} pag 99)
\end{itemize}




%\cite{digmanPhasorApproachFluorescence2008} per la somma di S e G

%\cite{claytonGraphicalRepresentationMulticomponent2004}

% citare il teorema di convoluzione \cite[Sez. 2.9.6]{oppenheim_discrete_2009}

%%%%%%%%%%%%%%%%%%%%%%%%%%%%%%%%%%%%%%%%%%%%%%%%%%%%%%%%%%%%%%%%%%%%%%%%%%%%%%%%%%%%
%%%%%%%%%%%%%%%%%%%%%%%%%%%%%%%%%%%%%%%%%%%%%%%%%%%%%%%%%%%%%%%%%%%%%%%%%%%%%%%%%%%%

\subsection{Funzione di risposta dello strumento} 

Lo studio dell'\textit{IRF} (\textit{Instrument Response Function}) è una parte fondamentale dell'analisi fasoriale. \\
L'emissione di fluorescenza è ottenuta grazie a un laser a impulsi periodici che produce un'eccitazione nelle molecole del fluoroforo presente nel campione. I fotoni emessi dallo stesso vengono catturati dalla scheda di acquisizione di un fotorivelatore; tuttavia questo non è in grado di distinguere tra i fotoni emessi dal campione e quelli dal laser.\\ % controllare l'ultima frase, potrebbe esserci un errore

% Il processo di separazione del segnale proveniente dalle due diverse fonti (campione e laser) è descritto da \textit{Ranjit et al.} \cite{ranjitFitfreeAnalysisFluorescence2018} e viene di seguito riportato.

Nonostante gli impulsi del laser siano molto brevi, la risposta del rivelatore introduce spesso un errore temporale, che in molti casi, può raggiungere ordini di grandezza paragonabili a quello dei tempi di vita di fluorescenza, su cui è incentrato lo studio. Si rivela dunque necessario separare i due segnali tramite una deconvoluzione, processo che richiede di conoscere la forma di entrambe le funzioni \cite{ranjitFitfreeAnalysisFluorescence2018}.\\

Il processo di acquisizione dell'\textit{IRF} presenta numerose criticità: in prima istanza, l'impossibilità di utilizzare una soluzione di scattering nel processo \textit{FLIM} a causa dei filtri ottici che bloccano i segnali con lunghezza d'onda pari a quella incidente.

In aggiunta a questo, il segnale subisce una distorsione (di solito un allargamento) nel passaggio attraverso l'obiettivo, la quale deve essere adeguatamente quantificata \cite{ranjitFitfreeAnalysisFluorescence2018}. \\

Emerge dunque un vantaggio significativo dell'analisi fasoriale rispetto al processo appena illustrato: non è necessario conoscere la forma precisa della \textit{IRF}, è sufficiente effettuare una calibrazione utilizzando un fluoroforo a decadimento mono-esponenziale a $\textit{\tau}$ noto.\\

Come ripreso nella sezione successiva, infatti,    % RICORDARSI DI MENZIONARLO PERÒ SENNÒ SEI UN CLOWN
la convoluzione nel dominio temporale si traduce come una moltiplicazione nello spazio di Fourier. La \textit{IRF} viene ridotta a un singolo punto sul piano dei fasori, che naturalmente non si troverà sul semicerchio universale in quanto non è un decadimento esponenziale \cite{ranjitFitfreeAnalysisFluorescence2018}. \\

Unendo queste due informazioni, è possibile ottenere il fasore del segnale totale come:

\begin{equation}
    z_{segnale} = z_{decadimento} \cdot z_{IRF}
    \label{eq:convoluzione_fasori}
\end{equation}

da cui è banale ricavare il fasore relativo al decadimento con un'operazione inversa.\\


%%%%%%%%%%%%%%%%%%%%%%%%%%%%%%%%%%%%%%%%%%%%%%%%%%%%%%%%%%%%%%%%%%%%%%%%%%%%%

% nella parte precedente della teoria devo spiegare bene come funziona l'additività delle fasi, ecc ecc
% qui volendo posso parlare per 6 ore


\subsubsection{Convoluzione tra IRF e segnale di fluorescenza}

È intereressante spiegare la ragione per cui le due funzioni (IRF e intensità di fluorescenza) si combinino come convoluzione invece che somma o moltiplicazione.\\

Se la IRF avesse la forma di una distribuzione \textit{Dirac-$\delta$}, tutte le molecole di fluoroforo verrebbero eccitate nello stesso istante, dunque avrebbero pari probabilità di decadimento al tempo \textit{t}, e si osserverebbe un decadimento di intensità di forma esponenziale decrescente \cite[Sez. 4.3.3]{mastersBookReviewPrinciples2008a}.\\

Si immagini di campionare l’impulso del laser in tanti istanti temporali successivi: all'istante $t_k$ l'intensità del laser sarà pari a $\textit{IRF}(t_k)$, mentre quella di fluorescenza seguirà la legge $I(t-t_k)$, dove l'argomento rappresenta il tempo trascorso dall'eccitazione.\\
Moltiplicando le due funzioni si ottiene un contributo dato da $\textit{IRF}(t_k)I(t-t_k)$.\\

Per ricavare l’intensità totale al tempo $t$, sarà quindi necessario sommare i contributi relativi a tutti gli istanti di eccitazione $t_k$, incluso quello coincidente con $t$ \cite[Sez. 4.3.3]{mastersBookReviewPrinciples2008a}:

\begin{equation}
    I_{tot}(t) = \sum_{t_k \le t} \, \textit{IRF}(t_k) \, I(t - t_k)
    \label{eq:convoluzione}
\end{equation}

La formula ottenuta si riconduce esattamente alla definizione matematica di convoluzione discreta tra una funzione $I$ e una funzione troncata \textit{IRF}.\\

%____________________________________________________________________________________

\paragraph{Osservazioni}
\begin{itemize}
    \item È giusto osservare che nell'implementazione nel codice dell'algoritmo di convoluzione non viene utilizzata la definizione appena ricavata, bensì una definizione operativa che fa uso della trasformata di Fourier. Per ulteriori approfondimenti si fa riferimento alla sezione \ref{sec:generazione_dei_segnali} (generazione dei segnali).
    \item Dato che l'acquisizione sperimentale avviene raggruppando i fotoni rilevati in \textit{bin} di larghezza pre-determinata $\Delta t$, sarebbe necessario pesare ogni contributo \textit{k-esimo} per questo valore. Nella definizione operativa questo termine viene inglobato dalla normalizzazione.\\
\end{itemize}


%%%%%%%%%%%%%%%%%%%%%%%%%%%%%%%%%%%%%%%%%%%%%%%%%%%%%%%%%%%%%%%%%%%%%%%%%%%%%%%%%%
%%%%%%%%%%%%%%%%%%%%%%%%%%%%%%%%%%%%%%%%%%%%%%%%%%%%%%%%%%%%%%%%%%%%%%%%%%%%%%%%%%


\section{Sviluppo del modello}\label{sec:sviluppo_del_modello} 


la stesura del codice si è divisa in due momenti fondamentali: la simulazione di segnali (e l'analisi dei suddetti)
e lo studio di immagini reali, acquisite tramite tecniche di microscopia a fluorescenza.\\

Sono state utilizzate le stesse funzioni per lo studio delle due tipologie di segnale.\\





\subsection{Generazione dei segnali}

% vedrò in seguito se spostare questa roba nella parte teorica

Come citato nella sezione (), %aggiungere riferimento ai cenni teorici nella sezione adeguata
i decadienti di fluorescenza seguono un andamento esponenziale decrescente;

\begin{equation}
    N(t) = N_0 e^{-\frac{t}{\tau}}
    \label{eq:exp_decay}
\end{equation}

dove \tau indica il tempo di vita di fluorescenza.\\

Si è dunque cercato di simulare un esperimento di fluorescenza così come avverrebbe in laboratorio: un fascio laser
periodico colpisce il campione causando un'eccitazione elettronica; al tempo \textit{t} un atomo ha una certa 
probabilità di trovarsi ancora allo stato eccitato. La dipendenza temporale di questa probabilità sarà 
dettata da \ref{eq:exp_decay}.\\

% Supponiamo ora di analizzare in tempo reale l'emissione di fluorescenza del campione per mezzo di un microscopio
% a fluorescenza. I fotoni emessi dal campione possono essere raggruppati in \textit{pixel}: l'intensità di ogni
% pixel avrà lo stesso andamento temporale di \ref{eq:exp_decay}.\\

Per ogni pulsazione del laser, si può misurare a quale \textit{t} dal momento dell'impulso avviene la 
diseccitazione (emissione di fluorescenza), e campionare i risultati ottenuti in \textit{t} di larghezza fissata
dalla scheda di acquisizione. Dopo un numero sufficiente di eventi, si sarà ottenuta una distribuzione esponenziale
decrescente, con tempo caratteristico \tau pari al tempo di vita di fluorescenza associato al fluoroforo.\\

Il fenomeno è stato simulato con una semplice generazione di distribuzione esponenziale, e il valore di \tau
verificato mediante il metodo della \textit{Binned Log Likelihood} (è stata utilizzata la libreria \textit{iMinuit}).

% reference teoria iminuit





\subsection{Deconvoluzione del segnale IRF}
\section{Test su immagini}\label{sec:test_su_immagini} 
\chapter{Conclusioni}\label{sec:conclusioni}



\printbibliography  % stampa bibliografia (duh)
% come citare: \cite{chiave} (es. Come derivato in \cite{dirac1928}...).

\end{document}


